\section{Introduction}
En una planta de proceso, un nodo sensor envía variables de baja velocidad ---temperatura, vibración y estado de una máquina--- hacia una estación de supervisión. El entorno contiene motores, contactores, variadores de frecuencia y cargas inductivas cuya conmutación produce perturbaciones electromagnéticas muy breves y de gran amplitud. Las mediciones en interiores identifican a los equipos con interruptores electromecánicos y a los motores como fuentes dominantes de ruido impulsivo de banda ancha \cite{blackard1993}, y su estadística se aparta del modelo gaussiano \cite{middleton1977}. En entornos industriales reales, este tipo de interferencia afecta enlaces de paquetes cortos y motiva técnicas de mitigación específicas \cite{zhan2022}.

El problema técnico no es aumentar la tasa de datos, sino conservar una señal reconocible y recuperar los bits cuando aparecen estos eventos. El trabajo se desarrolló en cuatro fases: modelado del canal con una perturbación impulsiva controlada; diseño de un transmisor con conformación de pulso y de un receptor digital propio; preparación del enlace sobre radios definidos por software (SDR), y caracterización del desempeño mediante curvas BER frente a $E_b/N_0$.

La Sección~\ref{sec:objetivos} define objetivos y alcance; la Sección~\ref{sec:metodologia} presenta el modelo, el transmisor, el receptor, la implementación SDR y la instrumentación; la Sección~\ref{sec:resultados} reporta y analiza los resultados, y la Sección~\ref{sec:conclusiones} resume las conclusiones.
